\documentclass[11pt]{amsart}

\usepackage[margin=1in]{geometry}
\usepackage{amsmath,amssymb,amsthm,mathtools}
\usepackage{hyperref}
\usepackage[shortlabels]{enumitem}
\usepackage{xcolor}
\definecolor{darkblue}{RGB}{0,0,139}

\newtheorem{theorem}{Theorem}[section]
\newtheorem{proposition}[theorem]{Proposition}
\newtheorem{lemma}[theorem]{Lemma}
\newtheorem{corollary}[theorem]{Corollary}

\theoremstyle{definition}
\newtheorem{definition}[theorem]{Definition}

\theoremstyle{remark}
\newtheorem{remark}[theorem]{Remark}

% Notation
\newcommand{\E}{\mathbb{E}}
\newcommand{\N}{\mathbb{N}}
\newcommand{\R}{\mathbb{R}}
\newcommand{\Prob}{\mathbb{P}}
\newcommand{\abs}[1]{\lvert #1 \rvert}
\newcommand{\norm}[1]{\lVert #1 \rVert}
\newcommand{\ip}[2]{\langle #1,#2\rangle}
\newcommand{\eps}{\varepsilon}
\newcommand{\sign}{\operatorname{sign}}
\newcommand{\card}[1]{\lvert #1 \rvert}
\newcommand{\tail}{\operatorname{tail}}
\newcommand{\restr}{\operatorname{head}}
\newcommand{\Span}{\operatorname{span}}
\newcommand{\supp}{\operatorname{supp}}
\DeclareMathOperator{\concFn}{Q}
\DeclareMathOperator{\smallBall}{\rho}

% Lean cross-references (breakable at _ and .; always upright mono)
\newcommand{\lnbrk}{\penalty100\hskip0pt plus 1pt\relax}
\newcommand{\leanfont}{\normalfont\ttfamily\color{darkblue}}
% block-level lean name (after a theorem statement)
\newcommand{\leanname}[1]{%
  \par\noindent\begingroup\footnotesize\leanfont
  \def\_{\char`\_\lnbrk}\def\.{\char`\.\lnbrk}%
  #1\endgroup\par\smallskip}
\newcommand{\privateleanname}[1]{%
  \par\noindent\begingroup\footnotesize\leanfont
  \def\_{\char`\_\lnbrk}\def\.{\char`\.\lnbrk}%
  #1\endgroup\par\smallskip}
% inline lean name (within running text)
\newcommand{\leanref}[1]{%
  {\footnotesize\leanfont
  \def\_{\char`\_\lnbrk}\def\.{\char`\.\lnbrk}%
  #1}}

\title{A Quantitative Proof of Stable Phase Retrieval}
\author{Pedro Abdalla Teixeira\textsuperscript{1}}
\author{Jaume de Dios Pont\textsuperscript{2}}
\author{Jo\~ao Pedro Ramos\textsuperscript{3}}
\author{Mitchell Taylor\textsuperscript{4}}

\begin{document}
\begin{abstract}
We give a fully explicit stability bound for phase retrieval in spans of
independent real-valued random variables under a uniform two-sided $L^1$
bound. If $\xi_1,\dots,\xi_n$ are independent, centered, unit-variance random
variables satisfying $A \le \E[\abs{\xi_i}] \le B$ for fixed $0 < A \le B < 1$,
then for any orthogonal unit vectors $a,b \in \R^n$ the modulus gap
$\sqrt{\E[(\abs{X_a}-\abs{X_b})^2]}$ is at least $c_\perp(A,B) > 0$, where
$c_\perp$ depends only on $A$ and $B$.

This result is the quantitative counterpart of the compactness proof of
stable phase retrieval \cite{compactness}. Each step of the compactness
argument---bounded probes, profile decomposition, head identification,
tail rigidity---is replaced by an explicit finite-dimensional estimate
with computable constants. The entire proof has been formalized in Lean~4
using the Mathlib library; Lean theorem names are recorded throughout.
\end{abstract}
\maketitle
\begingroup
\renewcommand{\thefootnote}{\arabic{footnote}}
\footnotetext[1]{Department of Mathematics, University of California, Irvine, Irvine, California 92697, USA}
\footnotetext[2]{Center for Data Science, New York University, New York, New York 10011, USA}
\footnotetext[3]{Instituto de Matem\'atica Pura e Aplicada (IMPA), Rio de Janeiro, RJ, Brazil}
\footnotetext[4]{Department of Mathematics, ETH Z\"urich, Z\"urich, Switzerland}
\endgroup

\section{Introduction}\label{sec:intro}

The phase retrieval problem asks whether a vector or function can be recovered
from phaseless measurements. In the real setting this means recovery up to
global sign. The finite-dimensional theory is by now extensive; see
\cite{bandeira2014saving, eldar2014stability} and the survey
\cite{grohs2020survey}. In infinite dimensions the situation is more
delicate: phase retrieval can hold in infinite-dimensional Hilbert spaces
\cite{cahill2016infinite}, but uniform stability may fail for continuous
frames \cite{alaifari2017continuous}, and natural measurement models such as
Gabor frames may be severely ill-posed \cite{alaifari2021gabor}.

We adopt the function-space point of view from \cite{FOTP}. Given a closed
subspace $E$ of a real function space $X$, we say that $E$ does \emph{stable
phase retrieval} with constant $C$ if
\[
\min_{\sigma \in \{\pm 1\}} \norm{f - \sigma g}_X
\le C \, \bigl\| \abs{f} - \abs{g} \bigr\|_X
\qquad
\text{for all } f, g \in E.
\]
The map $f \mapsto \abs{f}$ loses sign information, so the left-hand side
optimizes over a global sign flip; the inequality says that this is the only
obstruction to recovery. The main results of \cite{FOTP} show that many
order-continuous Banach lattices admit no infinite-dimensional subspace with
this property, while \cite{christ2022examples, FOTP} provide positive
examples in $L^2$-spaces. The orthogonal reduction of \cite{FOTP} is the
starting point of our argument: to verify stable phase retrieval on a
subspace, it suffices to check it for orthogonal pairs in the same
two-dimensional span.

In \cite{calderbank2022stable}, Calderbank, Daubechies, Freeman, and Freeman
proved stable phase retrieval for a model involving shifts and proposed the
following conjecture: if $(\xi_i)_{i \ge 1}$ is an orthonormal sequence of
independent mean-zero random variables in $L^2$, then
$\overline{\Span}(\xi_i)$ does stable phase retrieval if and only if there
exist $0 < A \le B < 1$ with
\[
A \le \norm{\xi_i}_{L^1} \le B
\]
for all but at most one $i$. The lower bound prevents $\xi_i$ from
concentrating near zero; the upper bound prevents it from being
Rademacher-valued (in which case $\abs{\xi_i}$ is constant and carries no
phase information). Together they ensure that each coordinate contributes
usable information about the signs of the coefficients.

Our main result proves the sufficiency direction of this conjecture with
a fully explicit stability constant.

\subsection{The main theorem}

Let $\xi_1, \dots, \xi_n$ be independent real-valued random variables on a
probability space $(\Omega, \mu)$, each satisfying
\[
\E[\xi_i] = 0,
\qquad
\E[\xi_i^2] = 1,
\qquad
A \le \E[\abs{\xi_i}] \le B
\]
for fixed $0 < A \le B < 1$. Given unit vectors $a, b \in \R^n$ with
$\ip{a}{b} = 0$, form the random linear combinations
$X_a = \sum_i a_i \xi_i$ and $X_b = \sum_i b_i \xi_i$. The \emph{modulus gap}
\[
\delta(a,b) = \sqrt{\E\bigl[(\abs{X_a} - \abs{X_b})^2\bigr]}
\]
measures how distinguishable $X_a$ and $X_b$ are from their absolute values
alone. After the orthogonal reduction, stable phase retrieval with constant $C$
is equivalent to $\delta(a,b) \ge \sqrt{2}/C$ for all orthogonal unit pairs.

\begin{theorem}[Quantitative orthogonal lower bound]\label{thm:main}
\leanname{Showcase\.quantitative\_orthogonal\_lower\_bound}
There exists a constant $c_\perp(A,B) > 0$, which can be chosen to be
\[
c_\perp(A,B)
= \tfrac{1}{2}
  \min\bigl(c_{\mathrm{cross}}(A),\;
  c_{\mathrm{line}}(A,B)\bigr),
\]
such that $c_\perp(A,B) \le \delta(a,b)$
for every pair of orthogonal unit vectors $a, b \in \R^n$.
The constant depends only on $A$ and $B$, not on the dimension $n$ or the
laws of $\xi_i$. The quantities $c_{\mathrm{cross}}$ and
$c_{\mathrm{line}}$ are defined in Section~\ref{sec:constants}.
\end{theorem}

\begin{remark}
The companion paper \cite{compactness} proves the same conclusion (without
explicit constants) by a compactness argument that also allows one
exceptional variable with no $L^1$ bound. The present proof makes each step
of that argument quantitative.
\end{remark}

\subsection{Proof strategy}\label{sec:strategy}

Rather than lower-bounding $\delta$ directly, we work with the auxiliary
quantity
\[
\eps(a,b) = \E\bigl[\abs{X_a^2 - X_b^2}\bigr],
\]
which is more amenable to algebraic manipulation. A Cauchy--Schwarz argument
(Theorem~\ref{thm:eps-le-modulus-app}) gives $\eps \le 2\delta$, so any lower
bound on $\eps$ transfers to $\delta$ at the cost of a factor of two.

The proof splits the coordinates into a \emph{head} (the few coordinates
where $a$ or $b$ has a large coefficient) and a \emph{tail} (the remaining
coordinates, each individually small). A sign dichotomy on the tail
(Lemma~\ref{lem:tail-dichotomy}) then determines which of two branches
applies:

\begin{enumerate}[label=(\roman*)]
\item \textbf{Concentrated branch} (Section~\ref{sec:head}).
For some sign $\tau = \pm 1$, the tail of $a - \tau b$ is small, so $a$ and
$\tau b$ nearly agree outside the head. Bounded probe functions extract
entries of the head difference matrix from the observable
$X_a^2 - X_b^2$, each entry being controlled by $K(A,B) \cdot \eps$.
The orthogonality $\ip{a}{b} = 0$ forces this matrix to have large Frobenius
norm, giving $\eps \ge c_{\mathrm{line}}(A,B)$.

\item \textbf{Diffuse branch} (Section~\ref{sec:blocks}).
Both $a + b$ and $a - b$ have substantial tail energy, spread across many
small coefficients. A greedy block partition, symmetrization via the product
space, Sperner's antichain bound, and a Chebyshev counting argument
together show that $X_a$ and $X_b$ are both unlikely to be near zero,
giving $\eps \ge c_{\mathrm{cross}}(A)$.
\end{enumerate}

\noindent
In both cases $\eps \ge \min(c_{\mathrm{cross}}, c_{\mathrm{line}})$, and
Theorem~\ref{thm:main} follows.


\section{Setup: the head--tail decomposition}\label{sec:setup}

Throughout, $a, b \in \R^n$ are unit vectors with $\ip{a}{b} = 0$, and
$\xi_1, \dots, \xi_n$ are independent random variables satisfying the moment
conditions of Section~\ref{sec:intro}. We write $X_c = \sum_i c_i \xi_i$ for
any coefficient vector $c \in \R^n$.

The first step is to separate the coordinates where $a$ and $b$ have large
coefficients from those where they are small. This is the quantitative
counterpart of the profile decomposition in \cite{compactness}: the large
coordinates form a finite set on which we can extract detailed information
via probes, while the small coordinates contribute a ``diffuse'' component
that is controlled by anticoncentration.

\begin{definition}[Head set, restriction, and tail]\label{def:head-set}
\leanname{Core\.FiniteModel\.headSetFin}
The \emph{head set} is
$H = H(A,a,b)
= \bigl\{i : \max(\abs{a_i}, \abs{b_i}) > \lambda(A)\bigr\}$.
For $S \subseteq \{1, \dots, n\}$ and $c \in \R^n$, write $\restr_S c$ for
the vector that agrees with $c$ on $S$ and vanishes elsewhere
(\leanref{restrictCoeff}), and
$\tail_S c = c - \restr_S c$ for its complement
(\leanref{tailCoeff}).
\end{definition}

Since each $i \in H$ contributes at least $\lambda(A)^2$ to
$a_i^2 + b_i^2 \le 2$, the head set is finite with
$\card{H} \le 2/\lambda(A)^2 \le d^*(A)$
(\leanref{headSetFin\_card\_le\_dStar}).
Conversely, every tail coordinate satisfies
$\max(\abs{(\tail_H a)_i}, \abs{(\tail_H b)_i}) \le \lambda(A)$
by construction
(\leanref{headSetFin\_tail\_max\_le\_lambda}).

\label{thm:head-card}%
\label{thm:tail-bound}%

\subsection{The head difference matrix}

The head difference matrix captures the discrepancy between $a$ and $b$ on
the large coordinates. It is a rank-two symmetric matrix whose Frobenius norm
is forced to be large by the orthogonality constraint whenever the tails
are small.

\begin{definition}[Difference matrix]\label{def:diff-matrix}
\leanname{Core\.FiniteModel\.finDH}
The head difference matrix $D_H(a,b)$ is the symmetric matrix with entries
\[
D_H(a,b)_{ij}
= (\restr_H a)_i (\restr_H a)_j
- (\restr_H b)_i (\restr_H b)_j
\]
and Frobenius norm
$\norm{D_H}_F = \sqrt{\sum_{i,j} D_H(a,b)_{ij}^2}$.
\end{definition}

\subsection{The epsilon--modulus bridge}

The two key quantities measuring phase-retrieval distinguishability are
\[
\eps(a,b) = \E\bigl[\abs{X_a^2 - X_b^2}\bigr],
\qquad
\delta(a,b)
= \sqrt{\E\bigl[(\abs{X_a} - \abs{X_b})^2\bigr]}.
\]
Since
$\abs{X_a^2 - X_b^2}
= \abs{X_a + X_b} \cdot \abs{X_a - X_b}$,
Cauchy--Schwarz and the variance identity
$\E[X_c^2] = \sum c_i^2$
(\leanref{Imported\.L2Theory\.randL2Sq\_finLinComb\_eq})
yield $\eps \le 2\delta$
(Theorem~\ref{thm:eps-le-modulus-app}). The rest of the
proof is devoted to showing $\eps$ is bounded below.


\subsection{Table of constants}\label{sec:constants}

All constants depend only on $A$ and $B$ and are explicitly computable.
The probe bound, gluing threshold, and head/tail cutoff are
\[
K = \max\!\bigl(\tfrac{4}{A^2},\,\tfrac{2}{1-B}\bigr),
\qquad
\theta = \tfrac{A^2}{8},
\qquad
\lambda = \tfrac{A^3}{2^{20}}.
\]
The head set has at most $d^* = 2 + 2^{41}/A^6$ elements, and the
diffuse threshold is $s_0 = A^4/2^{20}$. The two branch constants and
the stability constant are
\[
c_{\mathrm{cross}} = \tfrac{A^8}{2^{57}},
\qquad
c_{\mathrm{line}} = \tfrac{A^2}{2^8\, d^*\! K},
\qquad
c_\perp = \tfrac{1}{2}
\min(c_{\mathrm{cross}}, c_{\mathrm{line}}).
\]
The block partition in Section~\ref{sec:blocks} also uses
\[
M = \lceil 2^{18}/A^2 \rceil,
\quad
\eta(s) = \tfrac{s}{4M},
\quad
\lambda_0(s) = \tfrac{A\sqrt{s}}{2048},
\quad
r_0(s) = \tfrac{A^2\sqrt{s}}{65536}.
\]\label{sec:constants-end}

\section{The tail sign dichotomy}\label{sec:dichotomy}

The vectors $a + b$ and $a - b$ partition the energy of the pair $(a,b)$
into two complementary directions. On the tail $H^c$, at least one of these
directions must carry significant energy---but it may be that both do, or that
one is negligible. This dichotomy determines which branch of the proof
applies.

\begin{lemma}[Tail sign dichotomy]\label{lem:tail-dichotomy}
\privateleanname{Main\.UniformAB\.tail\_sign\_dichotomy}
Let $a, b \in \R^n$ be unit vectors with $\ip{a}{b} = 0$ and set
$H = H(A,a,b)$. Exactly one of the following holds:
\begin{enumerate}[label=(\roman*)]
\item \textbf{Diffuse:}
$\norm{\tail_H(a - b)}_2^2 \ge 16\, s_0(A)$ and
$\norm{\tail_H(a + b)}_2^2 \ge 16\, s_0(A)$.

\item \textbf{Concentrated:}
there exists $\tau \in \{+1, -1\}$ with
$\norm{\tail_H(a - \tau b)}_2 < \theta(A)/32$.
\end{enumerate}
\end{lemma}

\begin{proof}
If $\norm{\tail_H(a - b)}_2^2 < 16\, s_0$, take $\tau = 1$; if instead
$\norm{\tail_H(a + b)}_2^2 < 16\, s_0$, take $\tau = -1$. Since
$16\, s_0 = A^4/2^{16} = (\theta/32)^2$, the squared bound yields the norm
bound in~(ii). If neither holds, both squared norms are at least $16\, s_0$
and we are in~(i).
\end{proof}

\begin{remark}
In the concentrated case, $a$ and $\tau b$ nearly agree on the tail, so
their difference is concentrated on the finite head set $H$. In the diffuse
case, both $a \pm b$ have energy spread across many small tail coefficients,
precisely the regime where anticoncentration is effective.
\end{remark}

\section{Auxiliary estimates}\label{sec:auxiliary}

Before entering the two branches of the proof, we collect the three main
tools: bounded probe functions for reading off coefficients from the
quadratic observable $X_a^2 - X_b^2$; an $L^1$ lower bound for independent
sums; and a gluing estimate showing that if $a$ and $\tau b$ nearly agree
\emph{everywhere} (head and tail), then $\eps$ is already large.

\subsection{Probe functions}\label{sec:probes}

Given a centered random variable $\xi$ with $\E[\abs{\xi}] > 0$, one can
always find a bounded function of $\xi$ that acts as a dual to $\xi$ itself.
If additionally $\xi$ is not Rademacher-valued (i.e., $\E[\abs{\xi}] < 1$),
the same construction works for $\xi^2 - 1$. These are the probes that let
us read off individual entries of $D_H$ from $\eps$.

\begin{lemma}[Linear probe]\label{lem:linear-probe}
\leanname{Imported\.ElementaryBounds\.exists\_linear\_probe}
Let $\xi$ be a real random variable with $\E[\xi] = 0$ and
$\E[\abs{\xi}] \ge A > 0$. Then there exists a measurable
$g : \R \to \R$ with $\abs{g} \le 2/A$,\, $\E[g(\xi)] = 0$, and
$\E[\xi \, g(\xi)] = 1$.
\end{lemma}

\begin{proof}
Let $D = \E[\abs{\xi}]$ and $c = \E[\sign(\xi)]$, and set
$g(x) = (\sign(x) - c)/D$.
Then $\E[g(\xi)] = 0$ and
$\E[\xi \, g(\xi)] = \E[\abs{\xi}]/D = 1$. The bound $\abs{g} \le 2/D$
follows from $\abs{\sign(x) - c} \le 2$.
\end{proof}

\begin{lemma}[Quadratic probe]\label{lem:quadratic-probe}
\leanname{Imported\.ElementaryBounds\.exists\_quadratic\_probe}
Let $\xi$ be a real random variable with $\E[\xi] = 0$,
$\E[\xi^2] = 1$, and $\E[\abs{\xi}] \le B < 1$. Then there exists a
measurable $g : \R \to \R$ with $\abs{g} \le 2/(1-B)$,\,
$\E[g(\xi)] = 0$, and $\E[(\xi^2 - 1) \, g(\xi)] = 1$.
\end{lemma}

\begin{proof}
Apply the linear-probe construction to $\eta = \xi^2 - 1$, whose $L^1$ norm
satisfies
$\E[\abs{\xi^2 - 1}] \ge 1 - \E[\abs{\xi}] \ge 1 - B > 0$.
\end{proof}

\begin{remark}\label{rem:probe-extraction}
The linear probe recovers $c_i$ from $X_c = \sum_k c_k \xi_k$ via the
identity $\E[X_c \cdot g(\xi_i)] = c_i$: the centering $\E[g(\xi_i)] = 0$
kills the contribution of every $\xi_k$ with $k \ne i$, and the
normalization $\E[\xi_i g(\xi_i)] = 1$ picks out $c_i$. Likewise, the
quadratic probe gives $\E[X_c^2 \cdot g(\xi_i)] = c_i^2$.
\end{remark}

\subsection{The $L^1$ lower bound}\label{sec:l1-lower}

Independent sums of variables with positive $L^1$ norms cannot concentrate
near zero. This is the engine behind both the gluing estimate and the
diffuse branch: it converts $\ell^2$ energy of the coefficients into $L^1$
mass of the random sum.

\begin{theorem}[$L^1$ lower bound for independent sums]%
\label{thm:l1-lower}
\leanname{Imported\.ElementaryBounds\.finite\_indep\_sum\_l1\_lower}
Let $\zeta_1, \dots, \zeta_n$ be independent, mean-zero, unit-variance
random variables with $\E[\abs{\zeta_i}] \ge A > 0$ for all $i$. Then
for every $(c_1, \dots, c_n) \in \R^n$,
\[
\E\Bigl[\Bigl\lvert \sum_{i=1}^n c_i \zeta_i
\Bigr\rvert\Bigr]
\ge \frac{A}{4}
\sqrt{\sum_{i=1}^n c_i^2}\,.
\]
\end{theorem}

\subsection{The four-term gluing estimate}\label{sec:gluing}

Orthogonality forces $\norm{a + \tau b}_2 = \norm{a - \tau b}_2 = \sqrt{2}$.
If for some $\tau$ both the head of $a - \tau b$ and the tail of
$a + \tau b$ are small, then the complementary pieces---the tail of
$a - \tau b$ and the head of $a + \tau b$---must be close to $\sqrt{2}$.
The product $X_{a+\tau b} \cdot X_{a-\tau b} = X_a^2 - X_b^2$ then has a
large main term and small errors, yielding a lower bound on $\eps$.

\begin{lemma}[Four-term bound]\label{lem:four-term}
\privateleanname{Tail\.SignMismatch\.fin\_sign\_mismatch\_measure\_bound}
Let $a, b \in \R^n$ be unit vectors with $\ip{a}{b} = 0$, let
$H \subseteq \{1, \dots, n\}$, and let $\tau \in \{+1, -1\}$. Write
$c^+ = a + \tau b$ and $c^- = a - \tau b$. Then
\begin{align*}
\eps(a,b) \;\ge\;
&\bigl(\tfrac{A}{4}\bigr)^2\,
  \norm{\restr_H c^+}_2 \, \norm{\tail_H c^-}_2 \\
&{}- \norm{\restr_H c^+}_2 \, \norm{\restr_H c^-}_2
 - \norm{\tail_H c^+}_2 \, \norm{\tail_H c^-}_2 \\
&{}- \norm{\tail_H c^+}_2 \, \norm{\restr_H c^-}_2.
\end{align*}
\end{lemma}

\begin{proof}
The identity $X_a^2 - X_b^2 = X_{c^+} \cdot X_{c^-}$ holds since
$\tau^2 = 1$. Set $P = X_{\restr_H c^+}$, $X' = X_{\tail_H c^+}$,
$W = X_{\restr_H c^-}$, $V = X_{\tail_H c^-}$. Expanding
$(P + X')(W + V)$ and rearranging:
\[
\abs{PV}
\le \abs{X_{c^+} X_{c^-}} + \abs{PW} + \abs{X'V} + \abs{X'W}.
\]
Taking expectations, the left-hand side factors as
$\E[\abs{P}] \cdot \E[\abs{V}]$ by independence of disjoint supports
(\leanref{Imported\.L2Theory\.indepFun\_finLinComb\_disjoint}).
Each factor is at least $(A/4)$ times the $\ell^2$ norm of the
corresponding coefficient vector, by Theorem~\ref{thm:l1-lower}. The
three error terms are each bounded by the product of the relevant $L^2$
norms via Cauchy--Schwarz.
\end{proof}

\begin{proposition}[Gluing]\label{thm:gluing}
\leanname{Tail\.SignMismatch\.fin\_sign\_mismatch\_gluing}
Let $a, b \in \R^n$ be unit vectors with $\ip{a}{b} = 0$, let
$H \subseteq \{1, \dots, n\}$, and let $\tau \in \{+1, -1\}$.
If
\[
\norm{\restr_H(a - \tau b)}_2 \le \theta(A)/16
\quad \text{and} \quad
\norm{\tail_H(a + \tau b)}_2 \le \theta(A)/32,
\]
then $\eps(a,b) \ge \theta(A)/2$.
\end{proposition}

\begin{proof}
Write $c^+ = a + \tau b$ and $c^- = a - \tau b$. Set
$\alpha = \norm{\restr_H c^-}_2$ and $\beta = \norm{\tail_H c^+}_2$;
by hypothesis $\alpha \le \theta/16$ and $\beta \le \theta/32$.
From $\norm{c^\pm}_2^2 = 2$ we get
\[
\norm{\restr_H c^+}_2
= \sqrt{2 - \beta^2} \ge 1,
\qquad
\norm{\tail_H c^-}_2
= \sqrt{2 - \alpha^2} \ge 1.
\]
The main term in the four-term bound (Lemma~\ref{lem:four-term}) is
therefore at least
$(A/4)^2 \cdot 1 \cdot 1 = \theta/2$, while the three error terms sum to
at most $\alpha \sqrt{2} + \beta \sqrt{2} + \alpha\beta < \theta/4$.
The conclusion $\eps \ge \theta/2$ follows.
\end{proof}

\section{The concentrated branch}\label{sec:head}

In the concentrated case of the dichotomy
(Lemma~\ref{lem:tail-dichotomy}), there exists $\tau = \pm 1$ such that
$a$ and $\tau b$ nearly agree on the tail. Their difference is therefore
concentrated on the head set $H$, and the head difference matrix $D_H$
must be large. The strategy is to show that every entry of $D_H$ is
controlled by $K \cdot \eps$, while the Frobenius norm of $D_H$ is forced
to be at least $1$ by orthogonality. Together these give a lower bound on
$\eps$. If $\eps$ is too small, a bootstrapping argument via the gluing
estimate (Proposition~\ref{thm:gluing}) produces a contradiction.

\subsection{Entry bounds via probes}

\begin{proposition}[Head entry bound]\label{thm:head-entry}
\leanname{Head\.HeadMatrix\.head\_entry\_bound\_fin}
Let $a, b \in \R^n$ be unit vectors with $\ip{a}{b} = 0$ and let
$H = H(A,a,b)$. For all $i, j \in H$,
\[
\abs{D_H(a,b)_{ij}} \le K(A,B) \cdot \eps(a,b).
\]
\end{proposition}

\begin{proof}
\emph{Diagonal entries.}
Let $g$ be the quadratic probe for $\xi_i$
(Lemma~\ref{lem:quadratic-probe}). Writing
$X_c = c_i \xi_i + R$ with $R = \sum_{k \ne i} c_k \xi_k$ independent
of $\xi_i$, the computation in Remark~\ref{rem:probe-extraction} gives
$\E[X_c^2 \cdot g(\xi_i)] = c_i^2$.
Applying this to both $a$ and $b$:
\[
\abs{a_i^2 - b_i^2}
= \bigl\lvert\E\bigl[(X_a^2 - X_b^2) g(\xi_i)\bigr]\bigr\rvert
\le \frac{2}{1-B} \cdot \eps
\le K \cdot \eps.
\]

\emph{Off-diagonal entries.}
Let $g_i, g_j$ be linear probes for $\xi_i, \xi_j$
(Lemma~\ref{lem:linear-probe}). The same independence argument gives
$\E[X_c^2 \cdot g_i(\xi_i) g_j(\xi_j)] = 2 c_i c_j$, whence
\[
\abs{a_i a_j - b_i b_j}
\le \tfrac{1}{2} \cdot (2/A)^2 \cdot \eps
= \frac{2}{A^2}\, \eps
\le K \cdot \eps. \qedhere
\]
\end{proof}

\subsection{Frobenius control}

\begin{lemma}[Frobenius lower bound]\label{lem:head-frobenius}
\leanname{Head\.HeadMatrix\.head\_dominant\_lower\_bound\_fin}
Let $a, b \in \R^n$ be unit vectors with $\ip{a}{b} = 0$ and let
$H \subseteq \{1, \dots, n\}$. If
$\norm{\tail_H a}_2^2 + \norm{\tail_H b}_2^2 \le 1/100$,
then $\norm{D_H(a,b)}_F \ge 1$.
\end{lemma}

\begin{proof}
Write $\rho_a = \norm{\tail_H a}_2^2$ and
$\rho_b = \norm{\tail_H b}_2^2$, so that
$\norm{\restr_H a}_2^2 = 1 - \rho_a$ and
$\norm{\restr_H b}_2^2 = 1 - \rho_b$. The global orthogonality
$\ip{a}{b} = 0$ splits as
$\ip{\restr_H a}{\restr_H b}
= -\ip{\tail_H a}{\tail_H b}$,
and Cauchy--Schwarz gives
$\abs{\ip{\restr_H a}{\restr_H b}} \le \sqrt{\rho_a \rho_b}$.
Writing $z = \ip{\restr_H a}{\restr_H b}$, the Frobenius norm satisfies
\[
\norm{D_H}_F^2
= (1 - \rho_a)^2 + (1 - \rho_b)^2 - 2z^2.
\]
Since $\rho_a + \rho_b \le 1/100$ and
$z^2 \le \rho_a \rho_b \le (\rho_a + \rho_b)^2/4$,
one checks $\norm{D_H}_F^2 \ge 1$.
\end{proof}

Combining the entry bound with the Frobenius lower bound:
if $\norm{\tail_H a}_2^2 + \norm{\tail_H b}_2^2 \le 1/100$, then
$1 \le \norm{D_H}_F^2 \le \card{H}^2 (K\eps)^2$, giving
$\eps \ge 1/(\card{H} \cdot K)$
(\leanref{Head\.HeadMatrix\.threshold\_head\_lower\_bound\_fin}).

\subsection{The one-line branch}\label{sec:one-line}

We now close the concentrated case. The key algebraic ingredient is:

\begin{lemma}[Rank-one algebra]\label{lem:rank-one}
\privateleanname{Tail\.OneLine\.rank\_one\_sign\_selection}
For $x, y \in \R^m$, let $D(x,y)_{ij} = x_i x_j - y_i y_j$. Then
\[
\norm{x - y}_2^2 \cdot \norm{x + y}_2^2
\le 2\,\norm{D(x,y)}_F^2.
\]
\end{lemma}

\begin{proof}
Writing $s = \norm{x}_2^2$, $t = \norm{y}_2^2$, $c = \ip{x}{y}$, the
left-hand side is $(s + t)^2 - 4c^2$ and the right is
$2(s^2 + t^2 - 2c^2)$. The difference is $(s - t)^2 \ge 0$.
\end{proof}

\begin{lemma}[Head sign selection]\label{lem:head-sign-sel}
\privateleanname{Tail\.OneLine\.head\_sign\_selection}
Let $a, b \in \R^n$ be unit vectors with $\ip{a}{b} = 0$, let
$H \subseteq \{1, \dots, n\}$, and let $\tau \in \{+1, -1\}$.
If $\norm{\tail_H(a - \tau b)}_2 \le \theta(A)/32$ and every entry
of $D_H$ satisfies
$\abs{D_H(a,b)_{ij}} \le K(A,B) \cdot \eps(a,b)$,
then
\[
\norm{\restr_H(a + \tau b)}_2
\le 2\,\card{H} \cdot K(A,B) \cdot \eps(a,b).
\]
\end{lemma}

\begin{proof}
Apply Lemma~\ref{lem:rank-one} to
$x = \restr_H a$, $y = \tau\,\restr_H b$:
\[
\norm{\restr_H(a - \tau b)}_2^2
\cdot \norm{\restr_H(a + \tau b)}_2^2
\le 2\,\card{H}^2 (K\eps)^2.
\]
Pythagoras gives
$\norm{\restr_H(a - \tau b)}_2^2
= 2 - \norm{\tail_H(a - \tau b)}_2^2 \ge 1$.
Dividing through and taking square roots:
$\norm{\restr_H(a + \tau b)}_2
\le \sqrt{2}\,\card{H} \cdot K\eps
\le 2\,\card{H} \cdot K\eps$.
\end{proof}

\begin{proposition}[One-line branch]\label{thm:one-line}
\leanname{Tail\.OneLine\.one\_line\_branch}
Let $a, b \in \R^n$ be unit vectors with $\ip{a}{b} = 0$, let
$H \subseteq \{1, \dots, n\}$ with $\card{H} \le d^*(A)$, and let
$\tau \in \{+1, -1\}$. If
$\norm{\tail_H(a - \tau b)}_2 \le \theta(A)/32$,
then $\eps(a,b) \ge c_{\mathrm{line}}(A,B)$.
\end{proposition}

\begin{proof}
Suppose $\eps < c_{\mathrm{line}} = A^2/(2^8 d^* K)$. By
Proposition~\ref{thm:head-entry}, every entry of $D_H$ satisfies
$\abs{D_{ij}} \le K\eps$, so Lemma~\ref{lem:head-sign-sel} gives
\[
\norm{\restr_H(a + \tau b)}_2
< 2 d^* K \cdot \frac{A^2}{2^8 d^* K}
= \frac{A^2}{128} = \frac{\theta}{16}.
\]
Now swap the sign: the head of $a + \tau b$ has norm less than
$\theta/16$, and the tail of $a - \tau b$ has norm at most $\theta/32$.
The gluing estimate (Proposition~\ref{thm:gluing}) applied with sign
$-\tau$ gives $\eps \ge \theta/2 > c_{\mathrm{line}}$, a contradiction.
\end{proof}

\begin{corollary}\label{cor:headset-one-line}
\leanname{Tail\.OneLine\.headSet\_one\_line\_branch}
Let $a, b \in \R^n$ be unit vectors with $\ip{a}{b} = 0$ and set
$H = H(A,a,b)$. If
$\norm{\tail_H(a - \tau b)}_2 \le \theta(A)/32$
for some $\tau \in \{+1, -1\}$, then
$\eps(a,b) \ge c_{\mathrm{line}}(A,B)$.
\end{corollary}

\begin{proof}
Since $\card{H} \le d^*(A)$, Proposition~\ref{thm:one-line} applies.
\end{proof}

\section{The diffuse branch}\label{sec:blocks}

In the diffuse case, both $\tail_H(a \pm b)$ carry substantial $\ell^2$
energy, spread across many coordinates each of size at most $\lambda(A)$.
The goal is to show that $X_a$ and $X_b$ are both bounded away from zero
with positive probability, which forces
$\abs{X_a^2 - X_b^2}$ to be large.

The argument has three layers. First, we partition the tail coordinates into
blocks of comparable energy. Second, on the product space
$\Omega \times \Omega$, the symmetrized block sums satisfy a small-ball
bound via Sperner's antichain inequality. Third, a head--tail factorization
converts the anticoncentration into a lower bound on $\eps$.

\subsection{Greedy block partition}

\begin{lemma}[Block partition]\label{lem:block-partition}
\leanname{OneD\.Blocks\.exists\_block\_partition\_with\_bounds}
Let $(a_i)_{i=1}^n$ be real numbers satisfying $s > 0$,
$s \le \sum_{i=1}^n a_i^2$, and $a_i^2 \le \eta(A,s)$ for all $i$.
Then there exist pairwise disjoint
$G_1, \dots, G_{M(A)} \subseteq \{1, \dots, n\}$ with
\[
\eta(A,s) \le \sum_{i \in G_j} a_i^2 \le 2\,\eta(A,s)
\quad \text{for } j = 1, \dots, M(A).
\]
\end{lemma}

\begin{proof}
Build the blocks greedily. The total energy is at least
$s = 4M\eta$. At each step, when $m$ blocks remain, the unused energy
is at least $2m\eta$, so a nonempty qualifying subset exists. Choosing a
minimal qualifying subset $S$ with $\sum_{S} a_i^2 \ge \eta$: if this
sum exceeds $2\eta$, dropping any one element still leaves at least
$\eta$ (since each $a_i^2 \le \eta$), contradicting minimality.
\end{proof}

\subsection{Anticoncentration on the product space}%
\label{sec:finite-conc}

Fix a coefficient vector $c \in \R^n$ supported on $H^c$, and let
$G_1, \dots, G_M$ be the block partition from
Lemma~\ref{lem:block-partition}.
Pass to the product space $(\Omega \times \Omega, \mu \otimes \mu)$ and
define the differences
$D_i(\omega_1, \omega_2)
= \xi_i(\omega_1) - \xi_i(\omega_2)$,
which are independent and symmetric
(Theorem~\ref{thm:iIndepFun-differences-app}). The block sums
$Y_j = \sum_{i \in G_j} c_i D_i$ inherit both properties.

\begin{lemma}[Block mass]\label{lem:block-mass}
\privateleanname{OneD\.FiniteConcentration\.block\_two\_sided\_mass}
Let $\xi_1, \dots, \xi_n$ be independent, centered, unit-variance
random variables with $\E[\abs{\xi_i}] \ge A > 0$. Let
$(c_i)_{i \in G_j}$ satisfy
$\eta \le \sum_{G_j} c_i^2 \le 2\eta$, and define
$Y_j = \sum_{G_j} c_i D_i$ on $\Omega \times \Omega$.
Set $a_0 = (A/16)\sqrt{\eta}$. Then
\[
(\mu \otimes \mu)\bigl(\abs{Y_j} \ge a_0\bigr)
\ge A^2/256.
\]
\end{lemma}

\begin{proof}
Let $W = \sum_{G_j} c_i \xi_i$. The two-sided mass bound
(Theorem~\ref{thm:two-sided-mass-app}) gives
$\Prob(W \ge (A/16)\sqrt{\eta}) \ge A^2/256$ and likewise for $-W$.
Since $\Prob(W \ge 0) + \Prob(W \le 0) \ge 1$, combining the events
$\{W(\omega_1) \ge (A/16)\sqrt{\eta}\} \times \{W(\omega_2) \le 0\}$
and its mirror yields the stated bound on
$\{\abs{Y_j} \ge a_0\}$.
\end{proof}

\begin{lemma}[Small-ball bound]\label{lem:block-smallball}
\privateleanname{OneD\.FiniteConcentration\.smallBallMass\_sum\_symmetric\_le}
Let $Y_1, \dots, Y_M$ be independent symmetric random variables on
$\Omega \times \Omega$ with $M = M(A) \ge 256$. Suppose
$(\mu \otimes \mu)(\abs{Y_j} \ge a_0) \ge A^2/256$ for each $j$,
and let $0 < r < a_0$. Then for every $t \in \R$,
\[
(\mu \otimes \mu)
\Bigl(\abs{\textstyle\sum_{j=1}^{M} Y_j - t}
\le r\Bigr)
\le 1/8.
\]
\end{lemma}

\begin{proof}
Let $N = \#\{j : \abs{Y_j} \ge a_0\}$ and split:
\[
\Prob\Bigl(\abs{\textstyle\sum_j Y_j - t}
\le r\Bigr)
\le
\Prob\Bigl(\abs{\textstyle\sum_j Y_j - t}
\le r,\; N \ge 256\Bigr)
+ \Prob(N < 256).
\]
On $\{N \ge 256\}$, at least 256 of the $Y_j$ exceed $a_0 > r$ in
absolute value, so the Sperner sign-flip bound
(Theorem~\ref{thm:sperner-app}) gives a probability at most $1/16$.
For the second term: the expected count
$\sum_j \Prob(\abs{Y_j} \ge a_0) \ge M \cdot A^2/256 \ge 1024$,
so the Chebyshev bound (Theorem~\ref{thm:chebyshev-active-app}) gives
$\Prob(N < 256) \le 1/16$.
\end{proof}

\begin{proposition}[Symmetrized small-ball bound]\label{prop:symm-conc}
\leanname{OneD\.FiniteConcentration\.symmetrized\_small\_ball\_le}
Let $\xi_1, \dots, \xi_n$ be independent, centered, unit-variance
random variables with $A \le \E[\abs{\xi_i}] \le B$ for
$0 < A \le B < 1$. Let $c \in \R^n$ satisfy $s > 0$,
$s \le \sum c_i^2$, and $\abs{c_i} \le \lambda_0(A,s)$ for all $i$.
Then for every $t \in \R$,
\[
\mu\bigl(\abs{X_c - t} \le r_0(A,s)\bigr) \le 3/8,
\]
where $X_c = \sum_{i=1}^n c_i \xi_i$.
\end{proposition}

\begin{proof}
Unsymmetrization (Theorem~\ref{thm:unsymmetrization-app}) gives
\[
\mu\bigl(\abs{X_c - t} \le r_0\bigr)^2
\le
(\mu \otimes \mu)
\bigl(\abs{X_c(\omega_1) - X_c(\omega_2)} \le 2r_0\bigr).
\]
Write $X_c(\omega_1) - X_c(\omega_2) = \sum_j Y_j + T$, where $T$
is the residual outside the blocks. Since $\sum_j Y_j$ and $T$ act
on disjoint coordinates, the concentration function satisfies
$\concFn(\sum_j Y_j + T, 2r_0) \le \concFn(\sum_j Y_j, 2r_0)$
(\leanref{Imported\.L2Theory\.concFn\_add\_indep\_le}).
One checks that $r_0 < a_0$ and
$M(A) \ge 256$, so Lemma~\ref{lem:block-smallball} gives
$\concFn(\sum_j Y_j, 2r_0) \le 1/8$. Hence
$\mu(\abs{X_c - t} \le r_0) \le \sqrt{1/8} < 3/8$.
\end{proof}

\subsection{Joint anticoncentration and the cross estimate}

Applying Proposition~\ref{prop:symm-conc} to two coefficient vectors
simultaneously, we obtain:

\begin{corollary}[Cross anticoncentration]\label{cor:cross-anticonc}
\leanname{OneD\.Cross\.cross\_anticoncentration\_fin}
Under the hypotheses of Proposition~\ref{prop:symm-conc}, let
$c, d \in \R^n$ both satisfy $s > 0$, $s \le \sum c_i^2$,
$s \le \sum d_i^2$, and
$\abs{c_i}, \abs{d_i} \le \lambda_0(A,s)$. Then
\[
\Prob\bigl(r_0(A,s)
  < \min(\abs{X_c},\,\abs{X_d})\bigr)
\ge 1/4.
\]
\end{corollary}

\begin{proof}
At $t = 0$, Proposition~\ref{prop:symm-conc} gives
$\Prob(\abs{X_c} \le r_0) \le 3/8$ and
$\Prob(\abs{X_d} \le r_0) \le 3/8$.
The union bound gives $3/8 + 3/8 = 3/4$, so the complement has
probability at least $1/4$.
\end{proof}

It follows that the minimum of $\abs{X_c}$ and $\abs{X_d}$ has positive
second moment, uniformly over all shifts:

\begin{lemma}[Cross minimum-square bound]\label{lem:cross-min-sq}
\privateleanname{OneD\.Cross\.cross\_anticoncentration\_fin}
Let $c, d \in \R^n$ satisfy the hypotheses of
Corollary~\ref{cor:cross-anticonc} with $s = 8\,s_0(A)$.
Then for all $u, v \in \R$,
\[
\bigl(r_0(A,\, 8s_0)/4\bigr)^2
\le
\E\bigl[\min\bigl(\abs{X_c + u},\,
  \abs{X_d + v}\bigr)^2\bigr].
\]
\end{lemma}

\begin{proof}
Corollary~\ref{cor:cross-anticonc} applied with shifts $-u, -v$ gives
$\Prob(r_0 < \min(\abs{X_c+u}, \abs{X_d+v})) \ge 1/4$.
Writing $f = \min(\abs{X_c+u}, \abs{X_d+v})$, this means
$\E[f] \ge r_0/4$ (Markov), and $\E[f]^2 \le \E[f^2]$ (Jensen).
\end{proof}

The final ingredient converts anticoncentration of the tail into a
lower bound on $\eps$, by factoring $X_a^2 - X_b^2$ into head and tail
contributions:

\begin{lemma}[Head--tail factorization]\label{lem:head-tail-factor}
\privateleanname{OneD\.Cross\.finEps\_ge\_two\_cross\_min\_sq}
Let $a, b \in \R^n$ be unit vectors with $\ip{a}{b} = 0$ and let
$H \subseteq \{1, \dots, n\}$. Define tail-supported vectors
\[
c_i = \frac{a_i - b_i}{\sqrt{2}},\quad
d_i = \frac{a_i + b_i}{\sqrt{2}}
\quad \text{for } i \notin H,
\qquad
c_i = d_i = 0 \text{ for } i \in H.
\]
If there exists $r > 0$ with
$r \le \E[\min(\abs{X_c+u},\, \abs{X_d+v})^2]$
for all $u, v \in \R$, then $\eps(a,b) \ge 2r$.
\end{lemma}

\begin{proof}
The identity $X_a^2 - X_b^2 = X_{a+b} \cdot X_{a-b}$ and the
head--tail splitting give
\[
X_{a+b} \cdot X_{a-b}
= 2\,\bigl(X_d + h_+/\sqrt{2}\bigr)
    \bigl(X_c + h_-/\sqrt{2}\bigr),
\]
where $h_\pm = X_{\restr_H(a \pm b)}$ are the head sums. Thus
\[
\abs{X_a^2 - X_b^2}
\ge 2\,\min\bigl(\abs{X_c + h_-/\sqrt{2}},\,
\abs{X_d + h_+/\sqrt{2}}\bigr)^2.
\]
Since $X_c, X_d$ are supported on $H^c$ and $h_\pm$ on $H$, they
are independent
(\leanref{Imported\.L2Theory\.indepFun\_finLinComb\_disjoint}).
Conditioning on $h_\pm$ and applying the hypothesis
gives $\eps(a,b) \ge 2r$.
\end{proof}

\begin{proposition}[Diffuse cross branch]\label{prop:diffuse-cross}
\leanname{OneD\.Cross\.diffuse\_cross\_branch\_fin}
Let $a, b \in \R^n$ be unit vectors with $\ip{a}{b} = 0$ and
$H = H(A,a,b)$. If
\[
16\,s_0(A) \le \norm{\tail_H(a - b)}_2^2
\;\;\text{and}\;\;
16\,s_0(A) \le \norm{\tail_H(a + b)}_2^2,
\]
and $\max(\abs{(\tail_H a)_i}, \abs{(\tail_H b)_i})
\le \lambda(A)$ for all $i$, then
$\eps(a,b) \ge c_{\mathrm{cross}}(A)$.
\end{proposition}

\begin{proof}
The rotated coefficients $c, d$ from
Lemma~\ref{lem:head-tail-factor} satisfy
$\sum c_i^2 \ge 8s_0$, $\sum d_i^2 \ge 8s_0$, and
$c_i^2 \le 2\lambda^2 \le \lambda_0(A, 8s_0)^2$.
Lemma~\ref{lem:cross-min-sq} with $s = 8s_0$ gives
$\E[\min^2] \ge (r_0/4)^2$ for all shifts, and
Lemma~\ref{lem:head-tail-factor} converts this to
$\eps \ge 2(r_0/4)^2$. Unwinding the definitions:
\[
2\bigl(r_0(A,\, 8s_0)/4\bigr)^2
= \frac{A^8}{2^{55}}
\ge \frac{A^8}{2^{57}}
= c_{\mathrm{cross}}(A). \qedhere
\]
\end{proof}

\section{Proof of the main theorem}\label{sec:main}

\begin{theorem}[Epsilon lower bound]\label{thm:eps-lower}
\leanname{Main\.UniformAB\.orthogonal\_eps\_lower\_bound}
Let $a, b \in \R^n$ be unit vectors with $\ip{a}{b} = 0$, and let
$\xi_1, \dots, \xi_n$ be independent random variables with
$\E[\xi_i] = 0$, $\E[\xi_i^2] = 1$, and
$A \le \E[\abs{\xi_i}] \le B$ for $0 < A \le B < 1$. Then
\[
\eps(a,b)
\ge \min\bigl(c_{\mathrm{cross}}(A),\,
c_{\mathrm{line}}(A,B)\bigr).
\]
\end{theorem}

\begin{proof}
Set $H = H(A,a,b)$ and apply the dichotomy
(Lemma~\ref{lem:tail-dichotomy}).
In the diffuse case,
Proposition~\ref{prop:diffuse-cross} gives
$\eps \ge c_{\mathrm{cross}}$ (the tail coefficient bound follows from
the definition of $H$).
In the concentrated case, Corollary~\ref{cor:headset-one-line} gives
$\eps \ge c_{\mathrm{line}}$.
\end{proof}

\begin{proof}[Proof of Theorem~\ref{thm:main}]
Theorem~\ref{thm:eps-lower} gives
$\eps(a,b) \ge \min(c_{\mathrm{cross}}, c_{\mathrm{line}})$,
and the bridge inequality
(Theorem~\ref{thm:eps-le-modulus-app}) gives
$\eps(a,b) \le 2\,\delta(a,b)$.
Combining:
\[
\delta(a,b) \ge \frac{\eps(a,b)}{2}
\ge \frac{1}{2}
\min(c_{\mathrm{cross}}, c_{\mathrm{line}})
= c_\perp(A,B). \qedhere
\]
\end{proof}


\appendix
\section{Toolbox: Statements Used in the Main Proof}\label{app:toolbox}

This appendix collects the statements from the ``imported'' layer that are
invoked by the core proof (Sections~\ref{sec:head}--\ref{sec:main}). All
statements refer to a probability space $(\Omega, \mu)$ with independent
random variables $\xi_1, \dots, \xi_n$, and use the notation
$X_c = \sum_{i=1}^n c_i \xi_i$ for coefficient vectors $c \in \R^n$.
Proofs are available in the Lean formalisation.

\subsection{Symmetrization}

\begin{theorem}[Unsymmetrization]\label{thm:unsymmetrization-app}
\leanname{Imported\.Symmetrization\.unsymmetrization\_product\_space}
Let $W : \Omega \to \R$ be a square-integrable random variable on a
probability space $(\Omega, \mu)$. Then for every $r > 0$ and every
$t \in \R$,
\[
\mu\bigl(\abs{W - t} \le r\bigr)^2
\le
(\mu \otimes \mu)\bigl(\abs{W(\omega_1) - W(\omega_2)} \le 2r\bigr).
\]
\end{theorem}

\begin{theorem}[Differences preserve independence]%
\label{thm:iIndepFun-differences-app}
\leanname{Imported\.Symmetrization\.iIndepFun\_differences}
If $\zeta_1, \dots, \zeta_n$ are independent random variables on
$(\Omega, \mu)$, then the differences
$D_i(\omega_1, \omega_2) = \zeta_i(\omega_1) - \zeta_i(\omega_2)$ are
independent on $\Omega \times \Omega$.
\end{theorem}

\begin{theorem}[Desymmetrization $L^1$ bound]\label{thm:desymm-l1-app}
\leanname{Imported\.Symmetrization\.desymmetrization\_l1\_bound}
For any integrable random variable $W : \Omega \to \R$,
\[
\E_{\mu \otimes \mu}\bigl[\abs{W(\omega_1) - W(\omega_2)}\bigr]
\le
2 \, \E_\mu\bigl[\abs{W}\bigr].
\]
\end{theorem}

\subsection{$L^2$ Theory}

\begin{theorem}[Concentration drops under independent addition]%
\label{thm:concFn-add-app}
\leanname{Imported\.L2Theory\.concFn\_add\_indep\_le}
Let $U$ and $R$ be independent random variables on $(\Omega, \mu)$, and let
$r > 0$. Then
$\concFn_\mu(U + R, r) \le \concFn_\mu(U, r)$,
where $\concFn_\mu(W, r) = \sup_{t \in \R} \mu(\abs{W - t} \le r)$.
\end{theorem}

\begin{theorem}[$L^2$ norm of finite linear combinations]%
\label{thm:L2-finLinComb-app}
\leanname{Imported\.L2Theory\.randL2Sq\_finLinComb\_eq}
Let $\xi_1, \dots, \xi_n$ be independent, centered, unit-variance random
variables, and let $c = (c_1, \dots, c_n) \in \R^n$. Then
$\E\bigl[\abs{X_c}^2\bigr] = \sum_{i=1}^n c_i^2$.
\end{theorem}

\begin{theorem}[Epsilon--modulus inequality]\label{thm:eps-le-modulus-app}
\leanname{Imported\.L2Theory\.finEps\_le\_two\_mul\_finModulusGap}
Let $a, b \in \R^n$ with $\norm{a}_2 = \norm{b}_2 = 1$. Then
$\eps(a,b) \le 2 \, \delta(a,b)$.
\end{theorem}

\begin{theorem}[Jensen's inequality]\label{thm:jensen-app}
\leanname{Imported\.L2Theory\.sq\_integral\_le\_integral\_sq}
For any square-integrable random variable $W$ on $(\Omega, \mu)$,
$\bigl(\E[\abs{W}]\bigr)^2 \le \E[W^2]$.
\end{theorem}

\begin{theorem}[Independence of disjoint supports]%
\label{thm:indep-disjoint-app}
\leanname{Imported\.L2Theory\.indepFun\_finLinComb\_disjoint}
Let $c, d \in \R^n$ be coefficient vectors with
$\supp(c) \cap \supp(d) = \emptyset$, where
$\supp(c) = \{i : c_i \ne 0\}$. Then $X_c$ and $X_d$ are independent.
\end{theorem}

\subsection{Indicator Tails}

\begin{theorem}[Chebyshev indicator tail]\label{thm:chebyshev-indicator-app}
\leanname{Imported\.IndicatorTails\.chebyshev\_indicator\_tail}
Let $Z_1, \dots, Z_m$ be independent $\{0,1\}$-valued random variables with
$\sum_{j=1}^m \E[Z_j] \ge 1024$.
Then $\Prob\bigl(\sum_{j=1}^m Z_j < 256\bigr) \le 1/16$.
\end{theorem}

\begin{theorem}[Chebyshev active count]\label{thm:chebyshev-active-app}
\leanname{Imported\.IndicatorTails\.chebyshev\_active\_count}
Let $Y_1, \dots, Y_m$ be independent random variables, and let $\alpha > 0$.
If $\sum_{j=1}^m \Prob(\abs{Y_j} \ge \alpha) \ge 1024$, then
$\Prob\bigl(\#\{j : \abs{Y_j} \ge \alpha\} < 256\bigr) \le 1/16$.
\end{theorem}

\subsection{Elementary Bounds}

\begin{theorem}[Finite independent sum $L^1$ lower bound]%
\label{thm:finite-indep-l1-app}
\leanname{Imported\.ElementaryBounds\.finite\_indep\_sum\_l1\_lower}
Let $\zeta_1, \dots, \zeta_n$ be independent, mean-zero, unit-variance random
variables with $\E[\abs{\zeta_i}] \ge A > 0$ for all $i$. Then for every
$(c_1, \dots, c_n) \in \R^n$,
\[
\frac{A}{4} \sqrt{\sum_{i=1}^n c_i^2}
\;\le\;
\E\Bigl[\Bigl\lvert \sum_{i=1}^n c_i \zeta_i \Bigr\rvert\Bigr].
\]
\end{theorem}

\begin{theorem}[Two-sided mass of mean-zero sum]%
\label{thm:two-sided-mass-app}
\leanname{Imported\.ElementaryBounds\.two\_sided\_mass\_of\_mean\_zero}
Let $\zeta_1, \dots, \zeta_n$ be independent, mean-zero, unit-variance random
variables with $\E[\abs{\zeta_i}] \ge A > 0$ for all $i$, and let
$(c_1, \dots, c_n) \in \R^n$ with $\sigma = \sqrt{\sum c_i^2} > 0$. Then
\[
\Prob\Bigl(\tfrac{A}{16} \, \sigma
  \le \sum_{i=1}^n c_i \zeta_i\Bigr)
\ge A^2/256.
\]
\end{theorem}

\subsection{Rademacher: Sign Invariance}

\begin{theorem}[Sign-flip integral equality]\label{thm:sign-flip-app}
\leanname{Imported\.Rademacher\.SignInvariance\.sign\_flip\_integral\_eq}
Let $Z_1, \dots, Z_n$ be independent and symmetric random variables, and let
$(c_1, \dots, c_n) \in \R^n$. Then for any fixed sign pattern
$\varepsilon \in \{+1, -1\}^n$,
\[
\E\Bigl[\Bigl\lvert \sum_{i=1}^n c_i \varepsilon_i Z_i \Bigr\rvert\Bigr]
=
\E\Bigl[\Bigl\lvert \sum_{i=1}^n c_i Z_i \Bigr\rvert\Bigr].
\]
\end{theorem}

\subsection{Rademacher: Symmetric Sums}

\begin{theorem}[Symmetric Khintchine]\label{thm:sym-khintchine-app}
\leanname{Imported\.Rademacher\.SymmetricSums\.symmetric\_khintchine\_l1\_l2}
Let $Z_1, \dots, Z_n$ be independent and symmetric random variables with
$\E[Z_i^2] = 1$ and $\E[\abs{Z_i}] \ge A > 0$ for all $i$. Then for every
$(c_1, \dots, c_n) \in \R^n$,
\[
\frac{A}{2} \sqrt{\sum_{i=1}^n c_i^2}
\;\le\;
\E\Bigl[\Bigl\lvert \sum_{i=1}^n c_i Z_i \Bigr\rvert\Bigr].
\]
\end{theorem}

\subsection{Rademacher: Sperner}

\begin{theorem}[Sperner antichain bound]\label{thm:sperner-antichain-app}
\leanname{Imported\.Rademacher\.Sperner\.weighted\_sign\_interval\_antichain}
Let $w_1, \dots, w_n > 0$ with $w_i \ge \alpha > 0$ for all $i$, and let $I$
be an interval of diameter less than $2\alpha$. Then
\[
\#\Bigl\{\varepsilon \in \{+1,-1\}^n :
  \sum_{i=1}^n \varepsilon_i w_i \in I\Bigr\}
\;\le\;
\binom{n}{\lfloor n/2 \rfloor}.
\]
\end{theorem}

\subsection{Rademacher: Random Sperner}

\begin{theorem}[Sperner sign-flip]\label{thm:sperner-app}
\leanname{Imported\.Rademacher\.RandomSperner\.sperner\_sign\_flip\_bound}
Let $Y_1, \dots, Y_m$ be independent and symmetric random variables, let
$\alpha > 0$, $r > 0$ with $r < \alpha$, and let $m \ge 256$. Then for every
$t \in \R$,
\[
\mu\Bigl(
  \Bigl\{\abs{\textstyle\sum_{j=1}^m Y_j - t} \le r\Bigr\}
  \cap
  \Bigl\{\#\{j : \abs{Y_j} \ge \alpha\} \ge 256\Bigr\}
\Bigr)
\le 1/16.
\]
\end{theorem}


\begin{thebibliography}{10}

\bibitem{alaifari2017continuous}
R.~Alaifari and P.~Grohs.
\newblock Phase retrieval in the general setting of continuous frames for
  {B}anach spaces.
\newblock \emph{SIAM J.\ Math.\ Anal.}, 49(3):1895--1911, 2017.

\bibitem{alaifari2021gabor}
R.~Alaifari and P.~Grohs.
\newblock Gabor phase retrieval is severely ill-posed.
\newblock \emph{Appl.\ Comput.\ Harmon.\ Anal.}, 50:401--419, 2021.

\bibitem{compactness}
P.~Abdalla~Teixeira, J.~de~Dios~Pont, J.~P.~Ramos, and M.~Taylor.
\newblock A compactness proof of stable phase retrieval.
\newblock Preprint, 2025.

\bibitem{bandeira2014saving}
A.~S.~Bandeira, J.~Cahill, D.~G.~Mixon, and A.~A.~Nelson.
\newblock Saving phase: injectivity and stability for phase retrieval.
\newblock \emph{Appl.\ Comput.\ Harmon.\ Anal.}, 37(1):106--125, 2014.

\bibitem{cahill2016infinite}
J.~Cahill, P.~G.~Casazza, and I.~Daubechies.
\newblock Phase retrieval in infinite-dimensional {H}ilbert spaces.
\newblock \emph{Trans.\ Amer.\ Math.\ Soc.\ Ser.~B}, 3:63--76, 2016.

\bibitem{calderbank2022stable}
R.~Calderbank, I.~Daubechies, D.~Freeman, and N.~Freeman.
\newblock Stable phase retrieval for infinite dimensional subspaces of
  $L_2(\mathbb{R})$.
\newblock arXiv:2203.03135, 2022.

\bibitem{christ2022examples}
M.~Christ, B.~Pineau, and M.~A.~Taylor.
\newblock Examples of {H}\"older-stable phase retrieval.
\newblock \emph{Math.\ Res.\ Lett.}, 31(5):1339--1352, 2024.

\bibitem{eldar2014stability}
Y.~C.~Eldar and S.~Mendelson.
\newblock Phase retrieval: stability and recovery guarantees.
\newblock \emph{Appl.\ Comput.\ Harmon.\ Anal.}, 36(3):473--494, 2014.

\bibitem{FOTP}
D.~Freeman, T.~Oikhberg, B.~Pineau, and M.~A.~Taylor.
\newblock Stable phase retrieval in function spaces.
\newblock \emph{Math.\ Ann.}, 390(1):1--43, 2024.

\bibitem{grohs2020survey}
P.~Grohs, S.~Koppensteiner, and M.~Rathmair.
\newblock Phase retrieval: uniqueness and stability.
\newblock \emph{SIAM Rev.}, 62(2):301--350, 2020.

\end{thebibliography}

\end{document}
